\section{Background}

In this section, we recall some fundamental results related to the generalized Ricci flow equation. Given a smooth manifold, fix $g$ a Riemannian metric, \[H = \bigoplus_{k=1}^n H_k \in \Lambda^* T^* M,\] and a smooth function $f$. We recall the weighted sum defined in the introduction, and furthermore introduce
\begin{align*}
 \wwnorm{H}:= \sum_{k=1}^n\frac{k-1}k|H_k|^2, \qquad \wnorm{H} := \sum_{k=1}^n \frac{1}{k} \brs{H_k}^2.
\end{align*}
This data also determines notions of Ricci and scalar curvature:
\begin{defn} Given $(g, H, f)$ as above, the \emph{Ricci tensor} is
\begin{align*}
 \Rc^{H,f} := \Rc - \frac{1}{4} H^2 + \N^2 f - \frac{1}{2} \left( d^*_g H + i_{\N f} H \right) \in \Sym^2 T^* M \oplus \bigoplus_{k=0}^{n-1} \Lambda^k T^* M.
\end{align*}
Furthermore, the \emph{scalar curvature} is
\begin{align*}
 R^{H,f} = R - \frac{1}{4} \wnorm{H} + 2 \gD f - \brs{\N f}^2.
\end{align*}
\end{defn}

\begin{rmk} If a superscript in $\Rc^{H,f}$ or $R^{H,f}$ is dropped then the notation refers to the corresponding quantity with that term set to zero, i.e., $\Rc^f = \Rc + \frac{1}{2} \N^2 f$.
\end{rmk}

We note that in the case $H \in \Lambda^3$ and $f$ is constant, the Ricci tensor above is precisely the Ricci tensor of the unique metric-compatible connection with torsion $H$, which is a two-tensor with a symmetric and skew-symmetric part. For $H \in \Lambda^3$ and $f$ arbitrary, this tensor was defined in \cite{Streetsscalar} and named the twisted Bakry--\'Emery tensor. For general $H$ but constant $f$, this tensor is in the spirit of the generalized Ricci tensor used in \cite{FNS}. The coupling of the Ricci tensor to forms of arbitrary degree arises naturally in supergravity theories. Taking a hint from this, it may be possible to describe this Ricci curvature in general in terms of the curvature of a generalized connection on some augmented tangent bundle, as in the case of three-forms and the Bismut connection~\cite{GRFbook}.

For our purposes here, these definitions are justified by a key monotonicity formula for the scalar curvature along generalized Ricci flow. Given $(g_t, H_t, f_t)$ a solution to generalized Ricci flow as described in the introduction, we let
\begin{align*}
\square_f := \dt - \gD_f = \dt - \gD + \N f, \qquad \divg_f X := {\rm e}^{f} \divg \big({\rm e}^{-f} X\big)
\end{align*}
denote the forward weighted heat operator and weighted divergence. Before stating the result, we record some consequences of the fact that $H$ is closed which are left as exercises (cf.\ \cite[Lemma~3.19]{GRFbook} for the case $H$ is a three-form):
\begin{Lemma} \label{l:HBianchi} Given $(g, H)$ as above, one has
\begin{align*}
& \divg H^2 = - \IP{d^* H, H} + \frac{1}{2} d \wnorm{H},\qquad
 \divg \divg H^2 = \frac{1}{2} \gD \wnorm{H} + \sum_{k=1}^n \frac{1}{k} \IP{\gD_d H, H} + \brs{d^*_g H}^2,
\end{align*}
where for a $(k-1)$-form $\ga$ and $k$-form $\gb$, the notation $\IP{\ga, \gb}$ denotes the $1$-form uniquely defined by $\IP{\ga,\gb}(X) = \IP{\ga, i_X \gb}$.
\end{Lemma}

\begin{prop}[{\cite[Proposition 2.11]{Streetsscalar}}]\label{p:scalarmonotonicity} Given $(g_t, H_t, f_t)$ a solution to generalized Ricci flow, one has
\begin{align*}
\square_f R^{H,f} = 2 \big|\Rc^{H,f}\big|^2.
\end{align*}
\begin{proof} The result is claimed in \cite{Streetsscalar} without proof, so we include the short calculation here for convenience. Note furthermore that we are working here with the flow modified by diffeomorphisms generated by $\N f$. We compute the time derivative of each term in $R^{H,f}$ separately.
First, we compute that
	\begin{align}
 	\del_t R={}&-\IP{\Rc,\del_tg}+\divg\divg\del_t g-\Delta(\tr \del_t g)\nonumber\\
 ={}& 2\IP{\Rc,\Rc-\frac14 H^2}+\Delta R+\frac12\divg\divg H^2-\IP{\N R,\N f}-\frac12\Delta|H|^2, \label{eq: scalar}
 \end{align}
 	where we used the Bianchi identity and that
 	\begin{align*}
 	&\divg \N^2f= \N\Delta f+\Rc(\N f),\qquad\divg\divg \N^2f= \Delta^2 f+\frac12\IP{\N R,\N f}+\IP{\Rc,\N^2 f}.
 	\end{align*}
 	Then, we observe using Bochner's formula
 	\begin{align}
 	\del_t|\N f|^2={}& 2\left(\Rc+\N^2 f-\frac14 H^2\right)(\N f,\N f)+2\IP{\N f,\N\left( \Delta f-|\N f|^2 + \frac{1}{4} \wwnorm{H} \right)}\nonumber\\
 ={}& \Delta|\N f|^2-2|\N^2f|^2-\IP{\N f,\N |\N f|^2}+\frac12\big\langle\N\wwnorm{H},\N f\big\rangle\nonumber\\
 &-\frac12\IP{H^2,\N f\otimes \N f}.\label{eq: gradient}
 \end{align}
 Next, we compute 	
 	\begin{align}
 	\del_t \Delta f={}&\Delta \del_t f-\IP{\del_t g,\N^2 f}-\IP{\divg(\del_tg)-\frac12\N(\tr\del_tg),\N f}\nonumber\\
 ={}&\Delta^2 f-\Delta|\N f|^2+2\IP{\Rc+\N^2 f-\frac14 H^2,\N^2 f}-2\IP{-\divg\N^2f+\frac12\N\Delta f,\N f}\nonumber\\
 	& +\IP{-\frac12\divg H^2+\frac14\N |H|^2,\N f}+\frac14\Delta \wwnorm{H}.\label{eq: laplace}
 \end{align}
 	Finally, one has easily
 	\begin{align} \label{eq: h}
 	\del_t |H_k|^2=-k\IP{\del_t g, H_k^2}+2\IP{H_k,\Delta_d H_k-di_{\N f} H_k}.
 	\end{align}
 	Inserting \eqref{eq: scalar}, \eqref{eq: gradient}, \eqref{eq: laplace} and \eqref{eq: h} into the definition of $R^{H,f}$ yields
 \begin{align*}
 	\del_tR^{H,f}
 	={}&\Delta R^{H,f}+2\big|\Rc^{H,f}\big|^2-\big\langle\N R^{H,f},\N f\big\rangle,
 	\end{align*}
 where we used the identities for $\divg \N^2 f$ above and Lemma \ref{l:HBianchi}. The proposition follows.
\end{proof}
\end{prop}

\section{Wasserstein distance monotonicity for generalized Ricci flow}

Given a smooth connected manifold $M$, let $P(M)$ denote the space of Borel probability measures with finite second moments, i.e.,
\begin{align*}
 P(M):=\left\{\mu \text{ Borel probability measure}: \int_Md^2(x,x_0) {\rm d}\mu(x)<\infty \text{ for some }x_0\in M\right\}.
\end{align*} This space is naturally endowed with the Wasserstein distance $W$, defined for $\mu_1,\mu_2\in P(M)$ by the optimal transport problem
\begin{align*}
 W(\mu_1,\mu_2)^2:=\inf\int_{M\times M}d^2(x,y) {\rm d}\gamma(x,y),
\end{align*}
where the infimum is taken over all couplings $\gamma\in P(M\times M)$ with marginals $\gamma(\cdot\times M)=\mu_1$ and $\gamma(M\times \cdot)=\mu_2$. In the case $(M, g, {\rm e}^{-f} {\rm d}V)$ is a weighted Riemannian manifold, it is useful to consider the subspace $P^\infty(M)\subset P(M)$ consisting of smooth positive densities with respect to the weighted volume measure
\begin{align*}
 P^\infty(M):=\bigl\{\mu\in P(M)\colon {\rm d}\mu=\rho {\rm e}^{-f} {\rm d}V, \rho\in C^\infty(M),\ \rho>0\bigr\}.
\end{align*}
If $\mu\colon[0,1]\to P^\infty(M)$ is a smooth path, we write
\begin{align*}
 {\rm d}\mu(s)=\rho(s) {\rm e}^{-f} {\rm d}V
\end{align*}
and define $\phi(s)$ as a solution to the continuity equation
\begin{align}\label{eq: conteqstatic}
 \del_s\rho=-\divg_f(\rho\N\phi).
\end{align}
Such a $\phi(s)$ exists and is unique up to an additive constant.
Thus for such a path, we may define the Lagrangian
\begin{align*}
 E(\mu)=\frac12\int_0^1\int_M |\N\phi|^2 {\rm d}\mu {\rm d}s.
\end{align*}
A result known as the Benamou--Brenier formula shows that this formal notion of the length of a~path can be used to recover the Wasserstein distance, in the following sense (see \cite[Proposition~4.3]{OW05}):
\begin{thm}\label{thm: existence}
Let $\mu_1,\mu_2\in P^\infty(M)$ be probability measures. Then the infimum of $E$ over smooth curves in $P^\infty(M)$ satisfying the continuity equation and connecting these probability measures is $\frac12 W(\mu_1,\mu_2)^2$.
\end{thm}

In this section, we will analyze the monotonicity of the Wasserstein distance between two backward heat flows of probability measures under generalized Ricci flow. To begin, we record a~fundamental lemma, whose proof is elementary and left to the reader:

\begin{Lemma} \label{l:volumeev} Given $(g_t, H_t, f_t)$ a solution to generalized Ricci flow, one has
\begin{align*}
 \frac{{\rm d}}{{\rm d}t} {\rm e}^{-f} {\rm d}V = - R^{H,f} {\rm e}^{-f} {\rm d}V.
\end{align*}
\end{Lemma}

Also we derive a preliminary computation varying a certain integral along a curve of measures in a fixed time-slice.

